\documentclass[aspectratio=169]{beamer}
\usepackage[english]{babel}
\usepackage[utf8]{inputenc}
\usepackage[T1]{fontenc}
\usepackage{lmodern}
\usepackage{listings}
\usepackage{xcolor}
\usepackage{graphicx}
\usepackage{textcomp}
\DeclareUnicodeCharacter{2014}{---}
\DeclareUnicodeCharacter{2013}{--}
\DeclareUnicodeCharacter{2212}{-}
\DeclareUnicodeCharacter{2192}{\ensuremath{\rightarrow}}
\DeclareUnicodeCharacter{00B1}{\ensuremath{\pm}}
\DeclareUnicodeCharacter{00D7}{\ensuremath{\times}}
\DeclareUnicodeCharacter{00B7}{\ensuremath{\cdot}}
\DeclareUnicodeCharacter{20AC}{EUR}
\DeclareUnicodeCharacter{2070}{\textsuperscript{0}}
\DeclareUnicodeCharacter{00B9}{\textsuperscript{1}}
\DeclareUnicodeCharacter{00B2}{\textsuperscript{2}}
\DeclareUnicodeCharacter{00B3}{\textsuperscript{3}}
\DeclareUnicodeCharacter{2074}{\textsuperscript{4}}
\DeclareUnicodeCharacter{2075}{\textsuperscript{5}}
\DeclareUnicodeCharacter{2076}{\textsuperscript{6}}
\DeclareUnicodeCharacter{2077}{\textsuperscript{7}}
\DeclareUnicodeCharacter{2078}{\textsuperscript{8}}
\DeclareUnicodeCharacter{2079}{\textsuperscript{9}}
\DeclareUnicodeCharacter{2248}{\ensuremath{\approx}}
\DeclareUnicodeCharacter{2264}{\ensuremath{\le}}

% Colors for code
\definecolor{codebg}{RGB}{245,245,245}
\definecolor{codecomment}{RGB}{0,128,0}
\definecolor{codekeyword}{RGB}{0,0,255}
\definecolor{codestring}{RGB}{163,21,21}
\definecolor{bandcolor}{RGB}{200,50,50}

\lstdefinestyle{pythonstyle}{
    language=Python,
    backgroundcolor=\color{codebg},
    basicstyle=\ttfamily\scriptsize,
    keywordstyle=\color{codekeyword},
    commentstyle=\color{codecomment},
    stringstyle=\color{codestring},
    showstringspaces=false,
    breaklines=true,
    frame=single,
    rulecolor=\color{black},
    escapeinside={(*@}{@*)},
    moredelim=**[l][\color{bandcolor}\bfseries]{\#\ ---\ NEW},
    moredelim=**[l][\color{bandcolor}\bfseries]{\#\ ---\ CHANGED},
    literate=
      {á}{{\'a}}1 {é}{{\'e}}1 {í}{{\'i}}1 {ó}{{\'o}}1 {ú}{{\'u}}1
      {Á}{{\'A}}1 {É}{{\'E}}1 {Í}{{\'I}}1 {Ó}{{\'O}}1 {Ú}{{\'U}}1
      {ñ}{{\~n}}1 {Ñ}{{\~N}}1 {ü}{{\"u}}1 {Ü}{{\"U}}1
      {¿}{{?`}}1 {¡}{{!`}}1
      {—}{{---}}1 {–}{{--}}1 {−}{{-}}1
      {±}{{\ensuremath{\pm}}}1 {×}{{\ensuremath{\times}}}1
      {→}{{\ensuremath{\rightarrow}}}1
      {°}{{\textdegree}}1
      {·}{{\ensuremath{\cdot}}}1
      {€}{{\texteuro}}1
      {≈}{{\ensuremath{\approx}}}1
      {≤}{{\ensuremath{\le}}}1
      {⁰}{{\textsuperscript{0}}}1 {¹}{{\textsuperscript{1}}}1
      {²}{{\textsuperscript{2}}}1 {³}{{\textsuperscript{3}}}1
      {⁴}{{\textsuperscript{4}}}1 {⁵}{{\textsuperscript{5}}}1
      {⁶}{{\textsuperscript{6}}}1 {⁷}{{\textsuperscript{7}}}1
      {⁸}{{\textsuperscript{8}}}1 {⁹}{{\textsuperscript{9}}}1
}

\lstset{style=pythonstyle}

% Title slide macro: \lessontitle{Lesson Number}{Title}
\newcommand{\lessontitle}[2]{
    \title{Lesson #1: #2}
    \author{}
    \institute{Tec de Monterrey}
    \begin{frame}[plain]
        \titlepage
    \end{frame}
}

% Figure frame macro: \figureframe{Title}{Image path}
\newcommand{\figureframe}[2]{
    \begin{frame}{#1}
        \begin{center}
            \includegraphics[height=0.8\textheight,width=\textwidth,keepaspectratio]{#2}
        \end{center}
    \end{frame}
}


% Listings pulled from src/ with \lstinputlisting, so the slide is the file.
\lstset{
    basicstyle=\ttfamily\fontsize{8pt}{9.6pt}\selectfont,
    breaklines=true,
    breakatwhitespace=true,
    columns=fullflexible,
    keepspaces=true,
    xleftmargin=0.3em,
    framesep=2pt,
    aboveskip=0pt,
    belowskip=0pt
}

\newcommand{\resultframe}[3]{%
    \begin{frame}{#1}
        \begin{center}
            \includegraphics[height=0.62\textheight,width=\textwidth,keepaspectratio]{#2}
        \end{center}
        \vspace{0.25em}
        {\small #3\par}
    \end{frame}%
}

\newcommand{\claimframe}[2]{%
    \begin{frame}{#1}
        {\small #2\par}
    \end{frame}%
}

\date{}

\begin{document}

\lessontitle{07}{Tuning}

\section{This lesson}

\begin{frame}{This lesson}
Lesson 06 built the success-rate instrument and the evaluation counter. Today those tools are aimed at the parameters written at the top of every file so far.

\vspace{0.6em}
\begin{itemize}
    \item Tuning from single runs is reading noise.
    \item The honest currency is the margin over blind search at equal budget.
    \item A fixed setting is a compromise --- which is why Lesson 12 exists.
\end{itemize}
\end{frame}

% ================= tuning_01_single_runs =================
\begin{frame}{Paired evidence, not two independent rates}
The same seed index is used for both settings. For success indicators,
\[
d_i=I_{i,B}-I_{i,A}\in\{-1,0,1\},\qquad
\widehat\Delta=\bar d,\qquad
\widehat{\operatorname{SE}}(\widehat\Delta)=\frac{s_d}{\sqrt n}.
\]
This paired error replaces the independent-rate approximation. The final
36-cell map is exploratory for this landscape and seed set; it is not a
universal optimum or a confirmatory multiple-comparison study.
\end{frame}

\section{The book's method, faithfully}

\begin{frame}{The book's method, faithfully}
Chapter 7 is the book's laboratory chapter - 42 figures, three knobs
(crossover probability, mutation probability, population size), each turned
on its own landscape, each setting shown as one run, generation by
generation. It is the natural way to tune, and it is exactly the method
Lesson 06 demolished. This step keeps the book's protocol and its numbers -
population 16, mutation probability 0.2, crossover probabilities 0.0, 0.4,
0.7, seed 63 - on the course's landscape, and reads it the way the book
reads it. Then it reads it seven more times.
\end{frame}

\resultframe{The book's method, faithfully}{../figures/tuning_01_single_runs.png}{%
    One run per setting crowns pc = 0.7 by 0.0021; over 8 seeds the win counts go 0/3/5, the tightest split is 0.000004 (seed 3) and the widest 0.6284 (seed 1) — the ranking is unstable}

% ================= tuning_02_the_instrument =================
\section{The instrument}

\begin{frame}{The instrument}
Step 1 ran the book's method and watched it fail: one run per setting, a
ranking decided in the fourth decimal, win counts that flip with the seed.
Lesson 06 built the cure - a success rate over hundreds of runs, judged from
outside by a dense-grid reference and a stated tolerance. This step aims that
instrument at the first knob, crossover probability, on the course's
baseline: population 10, mutation probability 0.1, Lesson 01's sine
landscape, ten generations.
\end{frame}

\begin{frame}[fragile]{(1) the knobs become arguments}
\lstinputlisting[basicstyle=\ttfamily\fontsize{6pt}{7.2pt}\selectfont,firstline=113,lastline=142]{../src/tuning_02_the_instrument.py}
\end{frame}

\begin{frame}[fragile]{(2) standard\_error()}
\lstinputlisting[basicstyle=\ttfamily\fontsize{8pt}{9.6pt}\selectfont,firstline=159,lastline=167]{../src/tuning_02_the_instrument.py}
\end{frame}

\begin{frame}[fragile]{(3) measure()}
\lstinputlisting[basicstyle=\ttfamily\fontsize{7pt}{8.5pt}\selectfont,firstline=170,lastline=188]{../src/tuning_02_the_instrument.py}
\end{frame}

\begin{frame}[fragile]{(4) the course's baseline}
\lstinputlisting[basicstyle=\ttfamily\fontsize{8pt}{9.6pt}\selectfont,firstline=46,lastline=46]{../src/tuning_02_the_instrument.py}
\end{frame}

\resultframe{The instrument}{../figures/tuning_02_the_instrument.png}{%
    The observed crossover rate rises 55.6\% → 83.0\%; the 27.4-point endpoint gap is reported with the standard error of paired 0/1 differences. The table is still missing the cost column}

% ================= tuning_03_the_budget_knob =================
\section{The budget knob}

\begin{frame}{The budget knob}
Step 2 measured the crossover knob and found a clean monotone climb: 55.6\%
at 0.0 to 83.0\% at 1.0. It closed with the missing column. Crossover does
not come free - every crossed pair creates two new individuals, and every
new individual is a fitness evaluation, the only currency a GA spends. So
the dial that "improves" the algorithm also inflates its budget, and the
question becomes: how much of the climb is better search, and how much is
just more spending?
\end{frame}

\begin{frame}[fragile]{(1) the evaluation counter}
\lstinputlisting[basicstyle=\ttfamily\fontsize{7pt}{8.5pt}\selectfont,firstline=67,lastline=88]{../src/tuning_03_the_budget_knob.py}
\end{frame}

\begin{frame}[fragile]{(2) blind\_success()}
\lstinputlisting[basicstyle=\ttfamily\fontsize{8pt}{9.6pt}\selectfont,firstline=195,lastline=205]{../src/tuning_03_the_budget_knob.py}
\end{frame}

\begin{frame}[fragile]{(3) the margin column}
\lstinputlisting[basicstyle=\ttfamily\fontsize{8pt}{9.6pt}\selectfont,firstline=221,lastline=230]{../src/tuning_03_the_budget_knob.py}
\end{frame}

\resultframe{The budget knob}{../figures/tuning_03_the_budget_knob.png}{%
    Crossover 0.0 spends 20 evals/run, 1.0 spends 120; the margin over blind falls +30.6 → +0.8 points; the 100 extra evaluations buy the GA 27.4 points where 100 blind draws buy 57.2 — the knob is a budget knob}

% ================= tuning_04_mutation =================
\section{Mutation, and where the bell curve hides}

\begin{frame}{Mutation, and where the bell curve hides}
Every textbook draws the mutation curve as a bell: too little mutation and
the search stagnates, too much and it degenerates into random walk, with a
sweet spot in between. Gridin's Chapter 7 shows single runs at 0.0, 0.3 and
1.0 and gestures at the same shape. This step measures the knob properly -
seven settings, 500 runs each, crossover fixed at 0.8 - and the bell curve
is not there. The raw success rate climbs monotonically all the way to
mutation 1.0.
\end{frame}

\begin{frame}[fragile]{(1) the second knob}
\lstinputlisting[basicstyle=\ttfamily\fontsize{8pt}{9.6pt}\selectfont,firstline=199,lastline=204]{../src/tuning_04_mutation.py}
\end{frame}

\resultframe{Mutation, and where the bell curve hides}{../figures/tuning_04_mutation.png}{%
    The raw rate is monotone to mutation 1.0 (61.8\% → 89.8\%) — no bell curve; the margin bends: −10.9\% at 0.0 (worse than dice), peak +5.9\% near 0.2, −3.7\% at 1.0 — the bell curve lives in the margin}

% ================= tuning_05_population =================
\section{Population, the purest budget knob}

\begin{frame}{Population, the purest budget knob}
Crossover and mutation at least pretend to be about search strategy.
Population size pretends nothing: a bigger population is more evaluations,
full stop. It is the budget knob with no disguise, and Lesson 06 step 4
already caught it once - effectiveness climbs, efficiency falls. This step
re-measures it inside the tuning study, with the book's own settings
(6, 10, 20, 50, the book's mutation 0.2), to close the pattern before the
final map.
\end{frame}

\begin{frame}[fragile]{(1) the third knob}
\lstinputlisting[basicstyle=\ttfamily\fontsize{8pt}{9.6pt}\selectfont,firstline=200,lastline=206]{../src/tuning_05_population.py}
\end{frame}

\resultframe{Population, the purest budget knob}{../figures/tuning_05_population.png}{%
    Population 50 records 500/500 successes (95\% Wilson interval 99.2\%--100\%) at 550 evals/run; yield collapses 9.25 → 1.82 successes per 1000 evals; the margin peaks at population 10 (+5.9\%) and vanishes at 50 — "which is best?" has no answer without a currency}

% ================= tuning_06_the_grid =================
\section{The grid, and the verdict}

\begin{frame}{The grid, and the verdict}
Steps 2 to 5 turned one knob at a time and every knob confessed to being a
budget knob. The honest final question is whether they interact - whether
the best crossover depends on the mutation setting, the way the textbooks
assume when they recommend pairs like "0.8 and 0.1". This step measures the
full grid: six crossover settings times six mutation settings, population
10, 100 runs per cell, every cell priced against blind search at its own
budget.
\end{frame}

\begin{frame}[fragile]{(1) the grid}
\lstinputlisting[basicstyle=\ttfamily\fontsize{8pt}{9.6pt}\selectfont,firstline=192,lastline=204]{../src/tuning_06_the_grid.py}
\end{frame}

\begin{frame}[fragile]{(2) RUNS = 100}
\lstinputlisting[basicstyle=\ttfamily\fontsize{8pt}{9.6pt}\selectfont,firstline=47,lastline=47]{../src/tuning_06_the_grid.py}
\end{frame}

\resultframe{The grid, and the verdict}{../figures/tuning_06_the_grid.png}{%
    The (0, 0) corner recovers blind search (9.0\% ± 2.9\% vs 13.4\% for 10 blind draws — the instrument is calibrated); the best cell (1.0, 0.5) scores 90.0\% at 160 evals with a margin of −0.01\%; only 20 of 36 cells beat their own budget — the operators earn their keep in the cheap cells}

\section{Next}

\begin{frame}{Next}
Population is the budget knob with no disguise: a perfect success score can simply be bought, at a collapse in yield per evaluation.

\vspace{0.8em}
\textbf{Lesson 08 --- Black-Box Functions}

So far every landscape could be drawn and brute-forced. Lesson 08 takes the picture away. A fixed setting as a compromise is Lesson 12.
\end{frame}
\end{document}
